\section{Metodika výběrových šetření. Fáze dotazování, sestavení dotazníku, metody vážení}

\subsection*{Kvantitativní, kvalitativní a smíšený výzkum}
V dodaných materiálech je pro toto téma použitelný hlavně rozdíl mezi kvalitativním a kvantitativním výzkumem a část o sběru dat dotazníkem, rozhovorem a pozorováním.

\term{Kvalitativní výzkum} má za cíl porozumění, tvorbu nových hypotéz a nové teorie. Logika je induktivní: nejdříve se sbírají data, potom se v nich hledají pravidelnosti a významy. Používají se méně standardizované metody, například rozhovory, pozorování, terénní poznámky, audio a videozáznamy nebo deníky. Výhodou je hlubší pochopení kontextu a často vyšší validita. Nevýhodou bývá nižší reliabilita a horší zobecnitelnost.

\term{Kvantitativní výzkum} je založen na datech, proměnných, hypotézách a statistickém zpracování. Logika je deduktivní: z teorie nebo problému se vytvoří hypotézy, vyberou se proměnné, nasbírají se data a hypotézy se přijímají nebo zamítají. Výhodou je standardizace, relativně snadné porovnání odpovědí a možnost zobecnění výsledků na populaci. Nevýhodou může být redukce informace, protože respondent často vybírá jen z předem připravených kategorií.

\term{Smíšený výzkum} neboli metodologická triangulace kombinuje kvalitativní a kvantitativní metody. Například nejdříve rozhovory pomohou pochopit problém a formulovat hypotézy, potom dotazník ověří četnost jevů na větším vzorku.

\subsection*{Techniky shromažďování dat}
Materiály uvádějí několik technik, které lze využít i při výběrovém šetření.

\term{Analýza dokumentace} znamená zkoumání dostupných dokumentů, například výročních zpráv, směrnic, formulářů, popisů pracovních postupů, smluv, zápisů z jednání nebo projektové dokumentace. Umožňuje zjistit používanou terminologii, datové toky, formát dat, pravidla oběhu dokumentů a formální náležitosti. Rizikem je „paralýza analýzy“, tedy situace, kdy se sbírá příliš mnoho podkladů bez dostatečného přínosu.

\term{Pozorování} spočívá ve sledování činností v reálném prostředí. Je vhodné tam, kde lidé neumí přesně popsat, co skutečně dělají, nebo kde se chování liší od formálních postupů. U automatizovaných systémů lze podobnou roli plnit analýzou logů.

\term{Interview} je formální přímá komunikace tazatele s dotazovaným ve formě otázka - odpověď. Hodí se pro získávání obecných faktů, názorů, postojů a neformálních informací. Příprava rozhovoru zahrnuje definici cíle, seznámení s podklady, výběr respondenta, domluvu termínu a sestavení scénáře.

\term{Dotazník} je nepřímá a formalizovaná metoda sběru dat od většího počtu osob. Materiály jej popisují jako nejvíce formální a strukturovanou techniku. Dotazník se hodí pro zjišťování stanovisek, mínění, chování, charakteristik osob a empirických faktů. Výhodou je levnost, možnost oslovit více respondentů a snadnější kvantifikace. Nevýhodou je absence přímého kontaktu, nemožnost okamžitě vysvětlit nejasnou otázku a náročnost správného vyhodnocení.

\term{Anketa} je méně reprezentativní forma zjišťování názorů. Výsledky ankety nelze automaticky zobecnit na populaci, protože výběr respondentů nebývá reprezentativní.

\subsection*{Fáze dotazování}
Pro odpověď u zkoušky lze fáze dotazování sestavit z postupů uvedených v materiálech:
\begin{enumerate}
  \item \term{Vymezení cíle} - co chceme zjistit a proč. Bez jasného cíle nelze dobře formulovat otázky ani interpretovat odpovědi.
  \item \term{Vymezení cílové populace} - kdo má být předmětem šetření, například zákazníci, zaměstnanci, domácnosti nebo firmy.
  \item \term{Výběr respondentů} - kdo bude osloven. V materiálech není systematicky zpracovaná teorie výběru, ale pro praxi je nutné vědět, že výběr ovlivňuje reprezentativnost.
  \item \term{Návrh dotazníku nebo scénáře} - formulace otázek, jejich pořadí, typy odpovědí, škály a instrukce.
  \item \term{Pilotáž} - ověření, zda respondenti otázkám rozumějí a zda dotazník měří to, co měřit má.
  \item \term{Sběr odpovědí} - osobně, telefonicky, elektronicky nebo kombinovaně.
  \item \term{Kontrola a čištění dat} - kontrola úplnosti, validních hodnot, duplicit a logických nesrovnalostí.
  \item \term{Vyhodnocení} - popisné statistiky, porovnání skupin, případně modelování vztahů.
  \item \term{Interpretace a prezentace} - vysvětlení výsledků s ohledem na cíl šetření a omezení dat.
\end{enumerate}

\subsection*{Sestavení dotazníku}
Kvalitní dotazník musí být navržen tak, aby respondent rozuměl otázkám stejně jako autor dotazníku. Materiály zdůrazňují několik zásad.

Nejdříve je nutné jasně vysvětlit účel dotazníku a předem rozhodnout, jak budou odpovědi vyhodnocovány. Otázky mají být nezaujaté, jednoznačné, výstižné, přiměřeně krátké a terminologicky blízké respondentům. Je potřeba vyhnout se nápovědným otázkám, které respondenta tlačí k určité odpovědi, a dublovaným otázkám, které se ptají na dvě věci najednou.

\subsubsection*{Otevřené a uzavřené otázky}
\term{Otevřené otázky} dávají respondentovi prostor odpovědět vlastními slovy. Hodí se pro získání nových pohledů a kvalitativních informací. Nevýhodou je náročné vyhodnocení a obtížné porovnání mezi respondenty.

\term{Uzavřené otázky} mají předem dané možnosti odpovědi. Jsou rychlé, snadno porovnatelné a dobře kvantifikovatelné. Při jejich tvorbě je nutné zajistit, aby možnosti odpovědí byly srozumitelné, pokud možno vyčerpávající a vzájemně se nepřekrývaly. Pokud nejde pokrýt všechny varianty, lze použít možnost „jiné“ s doplněním.

\subsubsection*{Uspořádání otázek}
Materiály uvádějí tři přístupy:
\begin{itemize}
  \item \term{induktivní pyramida} - od detailních uzavřených otázek k obecnějším otevřeným;
  \item \term{deduktivní nálevka} - od obecných otevřených otázek k užším uzavřeným;
  \item \term{kosočtverec} - nejprve jednodušší uzavřené otázky, poté širší otevřené otázky a na závěr opět zúžení.
\end{itemize}

Pro dotazník je praktické sdružovat otázky tematicky, dávat důležité otázky spíše dříve a volit přehledný design formuláře.

\subsection*{Škálování odpovědí}
Škálování znamená přiřazení hodnot nebo symbolů vlastnostem, které chceme měřit. Materiály rozlišují:
\begin{itemize}
  \item \term{nominální měřítko} - pouze rozlišuje kategorie, například typ používaného programu nebo pohlaví;
  \item \term{ordinální měřítko} - umožňuje pořadí, například škála spokojenosti od velmi nespokojen po velmi spokojen;
  \item \term{intervalové měřítko} - vzdálenosti mezi hodnotami mají význam, ale nula nemusí znamenat neexistenci vlastnosti;
  \item \term{poměrové měřítko} - má absolutní nulu, lze porovnávat i poměry, například počet hodin nebo výše příjmu.
\end{itemize}

U zkoušky je dobré říct, že typ měřítka ovlivňuje, jaké statistické metody smíme použít. Ordinální odpovědi například nelze bez rozmyslu interpretovat jako plně metrické hodnoty.

\todohead
\begin{itemize}
  \item Doplnit pravděpodobnostní a nepravděpodobnostní výběr: náhodný, systematický, stratifikovaný, shlukový, kvótní, záměrný a snowball.
  \item Doplnit výběrovou chybu, nevýběrové chyby, non-response bias a pokrytí populace.
  \item Doplnit metody vážení: designové váhy, korekce na neodpovědi, post-stratifikace, kalibrace vah, normalizace vah.
  \item Doplnit praktický příklad výpočtu váhy respondenta.
\end{itemize}

\newpage
